% !TeX root = ../main.tex

\chapter{緒論}

\section{研究背景與動機}

隨著網際網路的蓬勃發展與行動裝置的普及，線上食譜平台（如 Food.com、Allrecipes 等）累積了數以百萬計的食譜內容與使用者互動紀錄。面對如此龐大的資訊量，使用者往往難以從海量的選項中找到真正符合自身口味偏好、飲食需求或烹飪條件的食譜。這種「資訊過載」(Information Overload) 問題催生了推薦系統 (Recommender System) 的廣泛應用與深入研究~\cite{su2009survey}。

傳統的推薦方法大多奠基於協同過濾 (Collaborative Filtering, CF) 技術~\cite{koren2009matrix}，透過分析使用者與物品之間的歷史互動矩陣來預測潛在偏好。然而，此類方法在面對資料稀疏性 (Data Sparsity) 與冷啟動 (Cold-start) 問題時，往往暴露出其架構上的根本限制——當使用者與物品的互動紀錄極度匱乏時，協同過濾模型便無法有效地捕捉潛在的偏好模式。

為突破上述瓶頸，學術界開始嘗試將知識圖 (Knowledge Graph, KG) 引入推薦系統~\cite{wang2018ripplenet, wang2019kgcn}。知識圖透過「頭實體——關係——尾實體 $(h, r, t)$」的三元組 (Triplet) 結構，構建出一個龐大且高度互連的語義網路。以食譜推薦為例，一道食譜可透過「包含食材」、「屬於菜系」、「適合場合」等多種關係，與其他食譜或屬性實體產生豐富的語義連結。這種結構化的背景知識能有效緩解資料稀疏性問題，並賦予推薦系統對食譜內涵的深層語義理解能力。

在此背景下，圖神經網路 (Graph Neural Networks, GNNs)~\cite{kipf2017gcn, velickovic2018gat} 成為融合知識圖與推薦系統的強力工具。其中，Wang 等人~\cite{wang2019kgat} 提出的知識圖注意力網路 (Knowledge Graph Attention Network, KGAT)，透過端到端的學習方式，將使用者與物品的互動矩陣和知識圖融合為一個協同知識圖 (Collaborative Knowledge Graph, CKG)，並運用注意力機制在多跳 (Multi-hop) 路徑上進行遞迴式的嵌入傳播，明確地模擬了高階連接性 (High-order Connectivity)，在推薦準確度上取得了突破性的進展。

然而，隨著模型架構複雜度的提升，圖神經網路推薦系統逐漸淪為難以解讀的「黑盒子」(Black-box)。使用者僅能看到推薦結果，卻無從得知系統為何做出此推薦。這種不透明性在食譜推薦場景中尤為關鍵——使用者可能存在過敏原限制、特殊飲食偏好或營養考量，若系統無法解釋推薦依據，將嚴重影響使用者的信任程度與採納意願~\cite{zhang2020explainable}。因此，如何在追求高準確度的同時，賦予推薦系統可解釋性 (Explainability)，使其能清楚闡述「為何推薦此食譜」，已成為當前研究的重要課題。

\section{現有方法之不足與研究缺口}

儘管 KGAT 在多個基準資料集上展現了優異的推薦效能，回顧現有文獻仍可辨識出以下三個尚未被充分探索的研究缺口：

\begin{enumerate}[leftmargin=2em]
  \item \textbf{領域適用性之空白：}KGAT 原論文~\cite{wang2019kgat}僅在 Amazon-Book、Yelp2018 及 Last-FM 三個電商或娛樂領域的基準資料集上進行驗證，尚未涉及食譜此種語義結構高度密集的領域。食譜知識圖具有與商品知識圖截然不同的圖拓撲特性——例如少數食材（如 \texttt{chicken}、\texttt{sugar}）往往連結數千筆食譜，形成高度偏斜的度分佈——這些特性可能對模型行為產生未預期的影響。

  \item \textbf{深度效應與模組貢獻分析仍有不足：}既有 KGAT 相關研究在報告實驗結果時，通常僅呈現最佳深度配置的效能數值，而較少系統性比較從 L=1 到 L=3 的效能變化趨勢。即使部分研究會討論注意力機制或知識圖的作用，其模組拆解多半侷限於有限設定，尚不足以完整釐清不同傳播深度下各元件的邊際貢獻如何變化。

  \item \textbf{深層 GNN 可解釋性驗證之匱乏：}多數可解釋性研究停留於定性案例展示，或僅在淺層模型（$L \leq 2$）上進行小樣本的忠實度 (Fidelity) 評估~\cite{yuan2022explainability}。隨著傳播深度增加，注意力路徑在多層非線性轉換後的 Fidelity 究竟如何演變，目前缺乏在深層 GNN（$L{=}3$）上的大規模 Fidelity 量化驗證。此外，傳統 Fidelity 指標在邊移除或保留過程中可能引發的分佈偏移 (Distribution Shift) 問題~\cite{agarwal2024fidelity}，亦鮮少在推薦場景中被討論。
\end{enumerate}

\section{研究目的}

基於上述研究缺口，本研究以 KGAT 為核心架構，在 Food.com 食譜推薦資料集上進行系統性的效能評估與可解釋性研究。具體而言，本研究旨在回答以下三個研究問題：

\begin{itemize}[leftmargin=2em]
  \item \textbf{RQ1：圖神經網路深度與知識圖整合如何影響推薦效能？} \\
  本研究將比較不同傳播深度下 KGAT 的效能變化，並在淺層設定下初步拆解注意力機制與知識圖的影響，並與 BPR-MF~\cite{rendle2009bpr}、LightGCN~\cite{he2020lightgcn}、NFM~\cite{he2017nfm} 等基準方法進行對比，以分析深度與知識整合在食譜推薦場景中的實際作用。
  \item \textbf{RQ2：不同傳播深度下，模型的訓練行為與泛化能力有何差異？} \\
    分析從 $L{=}1$ 到 $L{=}3$ 的深度變化對模型收斂速度、過擬合 (Overfitting) 現象及最終效能的影響。
  \item \textbf{RQ3：基於注意力機制的解釋路徑是否忠實地反映了模型的決策邏輯？} \\
    建構一套完整的可解釋性評估框架，以 500 位使用者為基礎，對三種不同深度的 KGAT 模型（$L{=}1$、$L{=}2$、$L{=}3$）進行跨模型 Fidelity 量化評估，驗證注意力路徑對模型預測的實際影響力。
\end{itemize}

\section{研究貢獻}

本研究的貢獻包含以下三點：

\begin{enumerate}[leftmargin=2em]
  \item \textbf{深度效應之實證分析：}本研究在 Food.com 食譜推薦場景下，比較了 $L{=}1$ 到 $L{=}3$ 的傳播深度對 KGAT 推薦效能的影響，呈現深度變化與效能提升之間的具體關聯，並支持高階連接性在此類知識圖推薦任務中的重要性。
  \item \textbf{淺層 KG 整合效果之觀察：}本研究在淺層設定 $L{=}1$ 下觀察到，引入知識圖未必必然帶來效能提升；相反地，在特定圖結構條件下，知識圖可能引入額外雜訊，進而稀釋原有的協同過濾訊號。此結果提醒我們，知識整合的效果可能與圖結構特性及傳播深度密切相關。
  \item \textbf{深層模型之 Fidelity 量化分析：}本研究建立從注意力路徑萃取到 Fidelity 量化的分析流程，並在 $L{=}1$、$L{=}2$、$L{=}3$ 三種深度的 KGAT 模型上，以 500 位使用者為基礎進行系統性的跨模型 Fidelity 比較評估。結果顯示，不同深度與路徑類型呈現不同的必要性與充分性分布，說明基於注意力的解釋雖具分析價值，但其解釋品質仍需結合量化指標審慎檢驗。
\end{enumerate}

\section{論文架構}

本論文共分為五章，各章節內容安排如下：

\textbf{第一章~緒論：}介紹研究背景、動機、前人缺口與三個核心研究問題。

\textbf{第二章~文獻探討：}回顧推薦系統的技術演進，包含協同過濾、圖神經網路、知識圖增強推薦及可解釋推薦等相關文獻，並以方法差異比較表歸納研究缺口。

\textbf{第三章~研究方法：}詳述本研究所採用的 KGAT 模型架構、知識圖建構流程、損失函數設計，以及基於 Fidelity 指標的可解釋性評估框架。

\textbf{第四章~實驗設計與結果分析：}呈現完整的實驗結果，包含主要效能對比、消融實驗分析及可解釋性評估結果，並對各項發現進行深入討論。

\textbf{第五章~結論與未來研究方向：}總結研究成果，回應三個研究問題，討論研究限制，並提出未來可改進的方向。
